%% kapa: exam=1 round=36 year=2025 period=2 form=A subject=accounting subject_ko=회계학 seq=8 number=48
%% source: Q-Net 감정평가사 기출문제 (cst00309) artlSeq=5231525 file=raw/1008/5231525/2025년 제36회 감정평가사 1차시험 2교시 문제지 원본.pdf
%% conversion: kordoc 4.24.51 -> tools/gen1.py (auto); review=proofread-sonnet; flags=none
\begin{question}{48}
(주)감평은 20×1년 초 사용 중인 기계장치(장부금액 ￦100,000, 공정가치 ￦40,000)를 (주)한국의 구축물(장부금액 ￦80,000, 공정가치 ￦70,000)과 교환하면서 (주)한국에 추가로 현금 ￦10,000을 지급하였다. 동 교환은 상업적 실질이 있으며, 기계장치의 공정가치가 구축물의 공정가치보다 명백하다. 한편 (주)감평은 교환으로 취득한 구축물의 내용연수와 잔존가치를 각각 4년과 ￦5,000으로 추정하였으며, 연수합계법으로 상각한다. (주)감평이 동 교환 시 인식할 유형자산처분손실(A)과 교환으로 취득한 구축물과 관련하여 20×1년에 인식할 감가상각비(B)는? (단, 교환으로 취득한 구축물은 영업에 사용하며, 원가모형을 적용한다.)
\par\noindent\hspace{3em}(A)\hspace{3.2em}(B)
\begin{choices}
  \item ￦30,000, ￦14,000
  \item ￦60,000, ￦14,000
  \item ￦60,000, ￦18,000
  \item ￦70,000, ￦14,000
  \item ￦70,000, ￦18,000
\end{choices}
\end{question}
\answer{3}
